\documentclass[12pt]{article}

\usepackage{sbc-template}
\usepackage[T1]{fontenc}
\usepackage[utf8]{inputenc}
\usepackage[brazil]{babel}
\usepackage{graphicx,url}
\usepackage{amsmath,amssymb,amsthm}
\usepackage{booktabs}
\usepackage{algorithm}
\usepackage{algpseudocode}
\usepackage{tikz}
\usetikzlibrary{fit}

\sloppy

% pseudocódigo em português
\floatname{algorithm}{Algoritmo}
\algrenewcommand\algorithmicrequire{\textbf{Entrada:}}
\algrenewcommand\algorithmicensure{\textbf{Saída:}}
\algrenewcommand\algorithmicwhile{\textbf{enquanto}}
\algrenewcommand\algorithmicdo{\textbf{faça}}
\algrenewcommand\algorithmicfor{\textbf{para}}
\algrenewcommand\algorithmicforall{\textbf{para cada}}
\algrenewcommand\algorithmicif{\textbf{se}}
\algrenewcommand\algorithmicthen{\textbf{então}}
\algrenewcommand\algorithmicelse{\textbf{senão}}
\algrenewcommand\algorithmicend{\textbf{fim}}
\algrenewcommand\algorithmicreturn{\textbf{devolva}}
\algrenewcommand\algorithmicfunction{\textbf{função}}
\algrenewcommand\algorithmicprocedure{\textbf{procedimento}}

\newtheorem{teorema}{Teorema}
\newtheorem{lema}{Lema}
\newtheorem{corolario}{Corolário}

\newcommand{\chiP}{\chi_P}
\newcommand{\LB}{\mathit{LB}}
\newcommand{\UB}{\mathit{UB}}

\title{Um Algoritmo \emph{Branch and Bound} para o\\
  Problema de Coloração Seletiva de Grafos}

\author{Carlos Alberto de Lima Neto, Guilherme Augusto Pondofe Santana}

\address{Departamento de Informática e Matemática Aplicada (DIMAp)\\
  Universidade Federal do Rio Grande do Norte (UFRN) -- Natal -- RN -- Brasil
}

\begin{document}

\maketitle

\begin{resumo}
  No Problema de Coloração Seletiva de Grafos, os vértices de um grafo estão
  divididos em $k$ grupos, e é preciso escolher exatamente um vértice de cada
  grupo e colorir os escolhidos com o menor número de cores. O problema é
  NP-difícil. Apresentamos um algoritmo exato do tipo \emph{branch and bound}
  com ramificação no estilo DSATUR, redução por dominância, limites inferiores
  por cliques e limite superior inicial dado pela heurística gulosa do nosso
  primeiro trabalho, melhorada por um guloso iterado. Provamos a correção do
  algoritmo, analisamos a sua complexidade e o aplicamos às 280 instâncias
  clássicas da literatura, com 60~s por instância. Ele provou o ótimo em 127
  instâncias e chegou ao valor de referência em 132 das 187 que têm
  referência. Em relação à heurística do primeiro trabalho, o gap médio caiu
  de 19\% e 43\% para no máximo 2,4\%. Somando os ótimos provados aos da
  literatura, a heurística sozinha atinge o ótimo em só 16 de 246 instâncias.
\end{resumo}

% ---------------------------------------------------------------------------
\section{Introdução}

Problemas de coloração de grafos modelam situações em que recursos precisam
ser atribuídos a elementos que entram em conflito. Na coloração seletiva, além
de colorir, é preciso escolher \emph{quais} elementos serão coloridos: os
vértices estão organizados em grupos e apenas um vértice de cada grupo entra na
solução. O problema foi proposto por \cite{li2000}, com o nome de
\emph{Partition Coloring Problem} (PCP), para modelar o roteamento e a
atribuição de comprimentos de onda em redes ópticas: cada grupo reúne as rotas
possíveis de uma conexão, duas rotas que compartilham um enlace são adjacentes
e cada cor é um comprimento de onda. Outras aplicações, como problemas de
escalonamento e de alocação, são discutidas por \cite{demange2015}.

Uma forma de ver o problema é a seguinte. Há tarefas divididas em categorias, e
de cada categoria basta executar uma tarefa. Duas tarefas são incompatíveis
quando acontecem no mesmo horário. Cada trabalhador executa um conjunto de
tarefas compatíveis entre si. Queremos escolher uma tarefa de cada categoria e
contratar o menor número de trabalhadores capaz de executá-las.

No primeiro trabalho, implementamos uma heurística gulosa que monta uma classe
de cor por vez. Ela é muito rápida, mas fica, em média, de 19\% a 43\% acima
dos valores ótimos publicados. Este trabalho apresenta um algoritmo
\emph{exato} para o problema, do tipo \emph{branch and bound}, e o compara com
essa heurística. As contribuições são: (i) um resumo do estado da arte em
algoritmos exatos para o PCP; (ii) um \emph{branch and bound} com ramificação
no estilo DSATUR, redução por dominância, dois limites inferiores, quebra de
simetria e limite superior inicial dado pela heurística do primeiro trabalho;
(iii) a prova de correção e a análise de complexidade do algoritmo; e (iv) a
aplicação do algoritmo às 280 instâncias clássicas da literatura, com
comparação à heurística do primeiro trabalho e aos valores ótimos publicados.

O restante do texto está organizado assim. A Seção~\ref{sec:problema} define o
problema e a Seção~\ref{sec:complexidade} discute a sua complexidade. A
Seção~\ref{sec:literatura} resume os algoritmos exatos já publicados e a
Seção~\ref{sec:tecnica} apresenta a técnica escolhida. A
Seção~\ref{sec:algoritmo} descreve o algoritmo implementado, cuja correção é
provada na Seção~\ref{sec:correcao} e cuja complexidade é analisada na
Seção~\ref{sec:analise}. A Seção~\ref{sec:experimentos} traz os experimentos e
a Seção~\ref{sec:conclusao} conclui o trabalho.

% ---------------------------------------------------------------------------
\section{Descrição do problema}\label{sec:problema}

Seja $G=(V,E)$ um grafo simples com $n=|V|$ vértices e $m=|E|$ arestas, e seja
$\mathcal{P}=\{V_1,\dots,V_k\}$ uma partição de $V$ em $k$ grupos não vazios.
Uma \emph{seleção} é um conjunto $S\subseteq V$ com $|S\cap V_i|=1$ para todo
$i$, isto é, um representante por grupo. Uma \emph{coloração própria} de $S$ é
uma função $f:S\to\{1,\dots,c\}$ tal que $f(u)\neq f(v)$ para toda aresta
$uv\in E$ com $u,v\in S$. O problema consiste em encontrar uma seleção $S$ e
uma coloração própria de $G[S]$ com o menor número possível de cores. Esse
mínimo é o \emph{número cromático de partição} de $G$, denotado por $\chiP(G)$:
\[
  \chiP(G)=\min\{\chi(G[S]) : S \text{ é uma seleção}\},
\]
onde $\chi$ é o número cromático usual. Arestas entre vértices de um mesmo grupo
nunca aparecem em $G[S]$ e podem ser ignoradas. Para um vértice $v\in V_i$,
escrevemos $N'(v)=N(v)\setminus V_i$ para o conjunto de vizinhos de $v$ fora
do seu grupo e $d_i(v)=|N'(v)|$. Também usamos $s=\max_i|V_i|$ para o tamanho
do maior grupo.

A Figura~\ref{fig:toy} mostra a instância de exemplo usada nos testes do
projeto. Os vértices ímpares formam um $K_4$, de modo que escolher o primeiro
vértice de cada grupo exige quatro cores. Escolhendo os vértices
$\{2,3,6,8\}$, duas cores bastam, e esse é o ótimo, pois nenhuma seleção é um
conjunto independente.

\begin{figure}[ht]
  \centering
  \begin{tikzpicture}[
      v/.style={circle,draw,minimum size=7mm,inner sep=0pt,font=\small},
      c1/.style={v,fill=black!20},
      c2/.style={v,fill=black,text=white},
      g/.style={draw,dashed,rounded corners,inner sep=2pt}]
    \node[v]  (1) at (0,1.6) {1};  \node[c1] (2) at (0,0) {2};
    \node[c1] (3) at (2.4,1.6) {3}; \node[v]  (4) at (2.4,0) {4};
    \node[v]  (5) at (4.8,1.6) {5}; \node[c1] (6) at (4.8,0) {6};
    \node[v]  (7) at (7.2,1.6) {7}; \node[c2] (8) at (7.2,0) {8};
    \foreach \a/\b in {1/3,3/5,5/7,2/4,6/8,2/5,4/7,1/6,3/8}
      \draw (\a) -- (\b);
    \draw (1) to[bend left=30] (5);
    \draw (3) to[bend left=30] (7);
    \draw (1) to[bend left=40] (7);
    \foreach \a/\b/\i in {1/2/1,3/4/2,5/6/3,7/8/4}
      \node[g,fit=(\a)(\b),label=below:{$V_\i$}] {};
  \end{tikzpicture}
  \caption{Instância de exemplo com $n=8$, $m=12$ e $k=4$. Os vértices
    preenchidos formam uma solução ótima: 2, 3 e 6 com a cor 1 (cinza) e 8 com
    a cor 2 (preto).}
  \label{fig:toy}
\end{figure}

% ---------------------------------------------------------------------------
\section{Complexidade do problema}\label{sec:complexidade}

A versão de decisão do problema (``existe seleção colorível com no máximo $c$
cores?'') pertence a NP: uma seleção e uma coloração formam um certificado
verificável em tempo $O(n+m)$. O problema é NP-difícil porque contém a coloração
de vértices como caso particular: quando todo grupo tem um único vértice, a
única seleção é $S=V$ e $\chiP(G)=\chi(G)$. Como decidir se $\chi(G)\le 3$ é
NP-completo \cite{karp1972,garey1979}, decidir se $\chiP(G)\le 3$ também é
\cite{li2000}. A coloração seletiva continua NP-difícil mesmo em classes de
grafos nas quais a coloração usual é fácil, conforme mostrado por
\cite{demange2014}.

A mesma redução vale para a aproximação. \cite{zuckerman2007} mostrou que, para
todo $\varepsilon>0$, não existe algoritmo polinomial que aproxime $\chi(G)$
dentro de um fator $n^{1-\varepsilon}$, a menos que P${}={}$NP. Na redução
acima, $k=n$, de modo que o PCP também não admite aproximação dentro de
$k^{1-\varepsilon}$ sob a mesma hipótese.

Assim, a menos que P${}={}$NP, nenhum algoritmo exato roda em tempo polinomial
no pior caso. O espaço de busca de um método ingênuo é grande: há
$\prod_{i=1}^{k}|V_i|$ seleções (até $2^k$ nas instâncias usadas, com grupos de
dois vértices) e, para cada uma, até $c^k$ funções de cor com $c$ cores. Uma
força bruta que enumera seleções e, para cada uma, testa
$c=1,2,\dots$ cores custa $O\big(\prod_i|V_i|\cdot k^{k}\cdot k^2\big)$ no pior
caso, inviável já para $k\approx 15$. Os algoritmos exatos práticos evitam
essa enumeração com limites e regras de poda, como descrito a seguir.

% ---------------------------------------------------------------------------
\section{Estado da arte em algoritmos exatos}\label{sec:literatura}

Os algoritmos exatos para o PCP derivam dos algoritmos exatos para a coloração
de vértices, e se dividem em duas famílias: os baseados em programação linear
inteira e os de enumeração implícita combinatória.

\paragraph{Formulação por representantes e \emph{branch-and-cut}.}
\cite{frota2010} adaptaram ao PCP a formulação por representantes da coloração
de vértices \cite{campelo2008}, na qual cada classe de cor é identificada por
um de seus vértices, o que elimina a simetria entre cores. Os autores
estudaram desigualdades válidas para o politopo dessa formulação e as usaram
em um algoritmo \emph{branch-and-cut}. Foi esse trabalho que introduziu os
conjuntos de instâncias random e ring usados até hoje.

\paragraph{Formulação de cobertura e \emph{branch-and-price}.} Na coloração
de vértices, a formulação de cobertura por conjuntos independentes, resolvida
por geração de colunas \cite{mehrotra1996}, fornece limites inferiores muito
fortes. \cite{hoshino2011} aplicaram ao PCP a decomposição de Dantzig-Wolfe da
formulação por representantes, obtendo um \emph{branch-and-price} em que cada
coluna é um conjunto independente que contém no máximo um vértice de cada
grupo. \cite{furini2018} propuseram um novo \emph{branch-and-price}, com
subproblema de conjunto independente de peso máximo, novas regras de
ramificação e uma busca adaptativa em vizinhança grande (ALNS) para gerar boas
soluções iniciais. Esse algoritmo resolveu pela primeira vez várias instâncias
em aberto e é o estado da arte nos conjuntos clássicos; os valores ótimos que
ele publica são a principal referência dos nossos experimentos.

\paragraph{Planos de corte em classes especiais.} \cite{seker2021}
propuseram um algoritmo exato de planos de corte para o caso em que o grafo é
perfeito, explorando o fato de que, nesse caso, o número cromático de cada
seleção é igual ao tamanho da sua maior clique. Eles também propuseram um
gerador de instâncias desse tipo.

\paragraph{\emph{Branch and bound} combinatório.} Para a coloração de
vértices, a principal alternativa à programação inteira são os algoritmos de
enumeração implícita derivados do DSATUR de \cite{brelaz1979}: a cada nó, o
vértice ainda sem cor com mais cores distintas na vizinhança recebe, em ramos
separados, cada cor viável ou uma cor nova. \cite{sewell1996} melhorou a regra
de desempate dessa escolha, e \cite{sansegundo2012} acrescentou limites
inferiores por cliques e outras melhorias, obtendo um dos algoritmos
combinatórios mais rápidos para a coloração exata. Como não dependem de
um resolvedor de programação linear, esses algoritmos podem ser implementados
em Python puro, como exige este trabalho. O algoritmo da
Seção~\ref{sec:algoritmo} adapta essa família ao PCP.

% ---------------------------------------------------------------------------
\section{Técnica escolhida: \emph{branch and bound}}\label{sec:tecnica}

O \emph{branch and bound} \cite{landdoig1960} resolve um problema de
otimização combinatória explorando uma árvore de subproblemas. Em cada nó, o
conjunto de soluções é dividido entre os filhos (a \emph{ramificação}), de
modo que toda solução do nó esteja em algum filho. O algoritmo guarda a melhor
solução já encontrada, chamada incumbente, cujo valor é o limite superior
$\UB$. Para cada nó, calcula-se um limite inferior $\LB$, um valor abaixo do
qual nenhuma solução do nó pode ficar. Em um problema de minimização, o nó é
podado quando o seu limite inferior é pelo menos $\UB$, pois ele não contém
solução melhor que a incumbente, ou quando não contém solução viável. Ao fim
da exploração, a incumbente é ótima: toda solução foi encontrada ou
descartada por um limite. Se a busca for interrompida antes,
o algoritmo ainda devolve uma solução viável ($\UB$) e o maior limite
inferior provado ($\LB$), e o valor ótimo fica no intervalo $[\LB,\UB]$.

A heurística do primeiro trabalho entra como limite superior inicial. Nas
instâncias de redes ópticas, há limites inferiores fáceis de calcular que
muitas vezes coincidem com o ótimo, e quando $\LB=\UB$ já na raiz nenhum nó
precisa ser explorado. A ramificação no estilo DSATUR fixa primeiro os grupos
mais restritos, de modo que a poda por inviabilidade age perto da raiz.

% ---------------------------------------------------------------------------
\section{Algoritmo implementado}\label{sec:algoritmo}

O algoritmo tem quatro etapas: redução por dominância, cálculo dos limites
inferiores na raiz, construção do limite superior inicial e busca em
profundidade. As subseções abaixo descrevem cada etapa, e o
Algoritmo~\ref{alg:bb} resume a busca.

\subsection{Redução por dominância}

Dizemos que $u$ \emph{domina} $v$ quando $u$ e $v$ estão no mesmo grupo e
$N'(u)\subseteq N'(v)$; se $N'(u)=N'(v)$, só o de menor índice domina o outro.
Um vértice dominado pode ser removido: trocar $v$ por $u$ em qualquer solução
não cria conflito (Lema~\ref{lem:dominancia}). Depois de remover os vértices
dominados de todos os grupos, as arestas que chegam a eles também deixam de
importar, o que pode criar novas dominâncias; por isso o processo é repetido
até não haver mudança. Nas instâncias nsfnet, por exemplo, uma rota cujos
enlaces estão todos contidos nos de outra rota da mesma conexão é removida.

\subsection{Limites inferiores na raiz}

\paragraph{Clique de grupos.} Seja $H$ o grafo cujos vértices são os grupos,
com $V_i$ e $V_j$ adjacentes em $H$ quando \emph{todo} vértice de $V_i$ é
adjacente a \emph{todo} vértice de $V_j$. Em qualquer seleção, os
representantes de dois grupos adjacentes em $H$ são vizinhos em $G$; logo, os
grupos de uma clique de $H$ precisam de cores distintas, e o tamanho $q$ de
uma clique de $H$ é um limite inferior. A clique é buscada por um
\emph{branch and bound} simples, limitado a $2\cdot 10^5$ nós; se o limite for
atingido, usamos a maior clique encontrada, que continua sendo um limite
válido. A clique também é usada na quebra de simetria
(Seção~\ref{sec:ramificacao}).

\paragraph{Cobertura por cliques.} Seja $\mathcal{F}=\{Q_1,\dots,Q_t\}$ uma
família de cliques de $G$ e $w(v)$ o número de cliques de $\mathcal{F}$ que
contêm $v$. Defina
\[
  W=\sum_{i=1}^{k}\ \min_{v\in V_i} w(v).
\]
Então $\chiP(G)\ge\lceil W/t\rceil$ (Lema~\ref{lem:cobertura}). Nas
instâncias de redes ópticas, cada enlace da rede induz uma clique (todas as
rotas que passam por ele), $w(v)$ é o comprimento da rota $v$ e o limite é o
clássico limite de carga: o total de enlaces usados dividido pelo número de
enlaces. Como o algoritmo não conhece a rede, as cliques são obtidas de forma
gulosa: enquanto existir uma aresta $vu$ ainda não coberta, criamos uma clique
maximal que contém $v$ e $u$. Em seguida, escolhemos a subfamília de duas
formas (acrescentando cliques enquanto $W/t$ cresce, ou retirando de todas as
cliques as que contribuem menos que a média $W/t$) e ficamos com o melhor
valor.

O limite inferior da raiz é $\LB=\max\{1,q,\lceil W/t\rceil\}$.

\subsection{Limite superior inicial}

A solução inicial é dada pela heurística gulosa do primeiro trabalho
(Algoritmo~\ref{alg:guloso}), que contrata um trabalhador (cor) por vez e dá a
ele o maior número possível de tarefas compatíveis, percorrendo os vértices em
ordem crescente de grau. Ela roda em $O(n\log n+cn+m)$.

\begin{algorithm}[ht]
  \caption{Guloso por classes de cor (Trabalho 1)}\label{alg:guloso}
  \begin{algorithmic}[1]
    \Require grafo $G=(V,E)$ e partição $\{V_1,\dots,V_k\}$
    \Ensure representante $r_i$ e cor $f_i$ de cada grupo $V_i$
    \State $L \gets$ vértices de $V$ em ordem crescente de $(\deg(v), v)$
    \State $P \gets \{1,\dots,k\}$ \Comment{grupos pendentes}
    \State $c \gets 0$
    \While{$P \neq \emptyset$}
      \State $c \gets c+1$; \quad $B \gets \emptyset$ \Comment{$B$: vizinhos da cor $c$}
      \ForAll{$v \in L$}
        \State $i \gets$ grupo de $v$
        \If{$i \in P$ \textbf{e} $v \notin B$}
          \State $r_i \gets v$; \quad $f_i \gets c$; \quad $P \gets P\setminus\{i\}$
          \State $B \gets B \cup N(v)$
        \EndIf
      \EndFor
    \EndWhile
    \State \Return $(r, f)$
  \end{algorithmic}
\end{algorithm}

Essa solução é melhorada pelo guloso iterado de \cite{culberson1996}, adaptado
à seleção. A cada iteração, as classes de cor da solução atual são
reordenadas (em ordem inversa, da maior para a menor ou aleatoriamente) e os
grupos são recoloridos com \emph{first-fit} nessa ordem: cada grupo recebe a
primeira cor em que algum dos seus vértices cabe, testando primeiro o
representante atual. Com isso, o número de cores nunca aumenta: os grupos da
$j$-ésima classe sempre cabem em uma das $j$ primeiras cores, pois cabiam
juntos na cor $j$ antes. O guloso iterado para quando atinge $\LB$, quando
passa 20.000 iterações seguidas sem reduzir o número de cores ou quando esgota
5~s (ou 10\% do limite de tempo, se for menor).

\subsection{Ramificação e quebra de simetria}\label{sec:ramificacao}

A busca trabalha sobre a instância reduzida. Um nó da árvore é uma
\emph{atribuição parcial}: alguns grupos já têm um par (representante, cor) e
os demais estão pendentes. Para cada cor $c$ já aberta, mantemos o conjunto
$B_c$ dos vizinhos dos vértices que estão na cor $c$; um vértice $v$ pode
receber a cor $c$ se, e somente se, $v\notin B_c$. Seja $u$ o número de cores
abertas e $\UB$ o valor da incumbente. Em cada nó:

\begin{enumerate}
  \item Para cada grupo pendente $V_i$, contamos as cores viáveis: as cores
    $c<u$ com algum vértice de $V_i$ fora de $B_c$, mais uma cor nova se
    $u+1<\UB$. Se algum grupo não tem nenhuma opção, o nó é podado.
  \item O grupo ramificado é o mais restrito, isto é, o com menos cores
    viáveis; no empate, fica o grupo com maior $\min_{v\in V_i}d_i(v)$, isto é,
    aquele cujo vértice menos conflitante ainda conflita com mais vértices.
    Um grupo com uma única opção é escolhido de imediato. Esse é o critério do
    DSATUR \cite{brelaz1979}, aplicado a grupos.
  \item Cada filho atribui ao grupo escolhido um par $(v,c)$ viável, com $v$
    em ordem crescente de grau. Cores ainda não usadas são equivalentes entre
    si, então só uma cor nova (a cor $u$) é tentada.
\end{enumerate}

A clique de grupos $Q=\{V_{g_0},\dots,V_{g_{q-1}}\}$ da raiz também quebra
simetria: o grupo $V_{g_j}$ só pode receber a cor $j$, e a busca começa com
$u=q$ cores abertas. Um ramo é podado quando $u\ge\UB$ (não pode melhorar a
incumbente) ou quando $\UB\le\LB$ (a incumbente já é ótima). Ao chegar a uma
folha (nenhum grupo pendente), a solução tem $u<\UB$ cores e passa a ser a
incumbente.

\begin{algorithm}[ht]
  \caption{Busca do \emph{branch and bound}}\label{alg:bb}
  \begin{algorithmic}[1]
    \Procedure{Busca}{$u$, $R$} \Comment{$u$: cores abertas; $R$: grupos pendentes}
      \If{$u \ge \UB$ \textbf{ou} $\UB \le \LB$} \Return \EndIf
      \If{$R=\emptyset$}
        \State incumbente $\gets$ atribuição atual; \quad $\UB \gets u$; \quad \Return
      \EndIf
      \ForAll{$V_i$ com $i\in R$}
        \State $D_i \gets \{j\}$ se $V_i=V_{g_j}\in Q$; senão $D_i\gets\{0,\dots,u-1\}$
        \State $C_i \gets \{c\in D_i : V_i\not\subseteq B_c\}$
        \If{$V_i\notin Q$ \textbf{e} $u+1<\UB$} $C_i\gets C_i\cup\{u\}$ \EndIf
        \If{$C_i=\emptyset$} \Return \Comment{poda por inviabilidade} \EndIf
      \EndFor
      \State $i \gets \arg\min_{i\in R}\,(|C_i|,\ -\min_{v\in V_i} d_i(v))$
      \ForAll{$c\in C_i$ \textbf{e} $v\in V_i$ com $v\notin B_c$}
        \State $r_i\gets v$; \quad $f_i\gets c$; \quad $B_c \gets B_c\cup N'(v)$
        \State \Call{Busca}{$\max(u, c+1)$, $R\setminus\{i\}$}
        \State desfaz a alteração em $B_c$
      \EndFor
    \EndProcedure
  \end{algorithmic}
\end{algorithm}

\subsection{Estratégias de busca}

O mesmo procedimento \textsc{Busca} é usado de duas formas.
\begin{itemize}
  \item \emph{Descendente}: \textsc{Busca} é chamada uma vez, com $\UB$
    igual ao valor da solução inicial; cada folha encontrada reduz $\UB$.
    É o \emph{branch and bound} clássico, bom para melhorar a solução.
  \item \emph{Ascendente}: para $t=\LB,\LB+1,\dots,\UB-1$, chamamos
    \textsc{Busca} com $\UB$ temporariamente igual a $t+1$, o que só aceita
    soluções com no máximo $t$ cores. Se a busca termina sem solução, está
    provado que $\chiP(G)>t$, e $\LB$ passa a $t+1$; se encontra uma
    solução, ela é ótima. É boa para provar limites inferiores.
\end{itemize}
A estratégia padrão, \emph{híbrida}, usa metade do tempo disponível na
busca descendente e, se ela não terminar, o restante na ascendente. Sem limite
de tempo, a busca descendente termina e a ascendente não é usada.

\subsection{Implementação}

O algoritmo foi implementado em Python~3, usando apenas a biblioteca padrão.
Conjuntos de vértices ($N'(v)$, $B_c$, os grupos) são representados por
inteiros usados como vetores de bits: o bit $v$ ligado indica que o vértice $v$
está no conjunto. Assim, testar $v\notin B_c$ ou contar os vértices de $V_i$
fora de $B_c$ são operações sobre inteiros, feitas em código nativo pelo
interpretador. O código está organizado assim:
\begin{itemize}
  \item \texttt{exato.py}: redução por dominância, limites inferiores, guloso
    iterado e o \emph{branch and bound};
  \item \texttt{sgcp.py}: leitura das instâncias, heurística do primeiro
    trabalho e validador, sem alterações em relação ao primeiro trabalho;
  \item \texttt{main.py}: interface de terminal. Um menu permite escolher o
    conjunto e depois a instância; a heurística e o \emph{branch and bound}
    são executados e um relatório com os dados da instância, o número de cores,
    o tempo e a validação de cada método, os limites inferior e superior, se o
    ótimo foi provado e a melhor solução (representante e cor de cada grupo e
    tarefas de cada trabalhador) com o seu valor é mostrado na tela e salvo em
    \texttt{resultados/};
  \item \texttt{experimentos.py}: executa os dois métodos em todas as
    instâncias e gera as tabelas de resultados;
  \item \texttt{testes.py}: testes automatizados, descritos na
    Seção~\ref{sec:experimentos}.
\end{itemize}
Toda solução passa pelo validador antes de ser mostrada ou salva. Para
executar, basta rodar, na pasta do projeto, \texttt{python main.py} (menu),
\texttt{python main.py instancias/exemplo/toy.txt} (instância direta, com
\texttt{--limite} para o tempo em segundos) ou \texttt{python experimentos.py}
(todas as instâncias). Em alguns sistemas o comando é \texttt{python3}.

% ---------------------------------------------------------------------------
\section{Prova de correção}\label{sec:correcao}

Escrevemos $\chiP$ para o valor ótimo da instância original.

\begin{lema}\label{lem:dominancia}
  A redução por dominância não altera o valor ótimo.
\end{lema}

\begin{proof}
  Remover vértices não cria soluções novas, então o ótimo não diminui. Para
  ver que ele não aumenta, seja $(S,f)$ uma solução e $v\in S\cap V_i$ um
  vértice dominado por $u\in V_i$, isto é, $N'(u)\subseteq N'(v)$. Troque $v$
  por $u$, mantendo a cor $f(v)$. Todo vizinho de $u$ na nova seleção pertence
  a $N'(u)\subseteq N'(v)$ e, portanto, já era vizinho de $v$, com cor
  diferente de $f(v)$. A nova solução é própria e usa o mesmo número de cores.
  Repetindo a troca, obtemos uma solução com o mesmo número de cores que só
  usa vértices não dominados (a relação de dominância é acíclica graças ao
  desempate pelo menor índice). Removido $v$, as arestas que chegam a ele não
  aparecem em nenhuma seleção da instância reduzida, então apagá-las preserva
  o conjunto de soluções; as rodadas seguintes aplicam o mesmo argumento à
  instância reduzida.
\end{proof}

\begin{lema}\label{lem:clique}
  Se $Q$ é uma clique de $H$ com $q$ grupos, então $\chiP\ge q$.
\end{lema}

\begin{proof}
  Para quaisquer dois grupos de $Q$, todo vértice de um é adjacente a todo
  vértice do outro, em particular os seus representantes. Logo, os $q$
  representantes formam uma clique em $G[S]$, para toda seleção $S$, e
  precisam de $q$ cores distintas.
\end{proof}

\begin{lema}\label{lem:cobertura}
  Para toda família de $t\ge1$ cliques de $G$, $\chiP\ge\lceil W/t\rceil$.
\end{lema}

\begin{proof}
  Seja $(S,f)$ uma solução ótima. Cada classe de cor $f^{-1}(c)$ é um
  conjunto independente, logo contém no máximo um vértice de cada clique
  $Q_j$; assim, $\sum_{v\in f^{-1}(c)}w(v)\le t$. Somando sobre as $\chiP$
  classes, $\sum_{v\in S}w(v)\le t\,\chiP$. Por outro lado, $S$ tem um vértice
  de cada grupo, então $\sum_{v\in S}w(v)\ge\sum_i\min_{v\in V_i}w(v)=W$.
  Daí $\chiP\ge W/t$ e, como $\chiP$ é inteiro, $\chiP\ge\lceil W/t\rceil$.
\end{proof}

\begin{lema}\label{lem:caminho}
  Considere uma chamada de \textsc{Busca} com valor $\UB$ e suponha que exista
  uma solução da instância reduzida com $c^\star<\UB$ cores. Então a busca
  encontra uma solução com no máximo $c^\star$ cores, a menos que antes disso
  $\UB$ caia para um valor $\le c^\star$ ou para $\LB$.
\end{lema}

\begin{proof}
  Seja $(S,f)$ essa solução. Como os grupos de $Q$ estão em uma clique, eles
  têm cores distintas em $f$; renomeando as cores, podemos supor que
  $V_{g_j}$ tem a cor $j$. Construímos um caminho da raiz até uma folha que
  ``segue'' $(S,f)$, mantendo um mapeamento injetivo $\pi$ das cores de $f$
  para as cores da busca, com $\pi(j)=j$ para $j<q$. Em cada nó, seja $V_i$ o
  grupo escolhido e $v$ o seu representante em $S$. Se $f(v)$ já está
  mapeada, tomamos o filho $(v,\pi(f(v)))$; caso contrário, $V_i\notin Q$ e
  tomamos o filho $(v,u)$, a cor nova, fazendo $\pi(f(v))=u$. Por
  construção, a cor $\pi(f(v))$ da busca só contém vértices de $S$ com cor
  $f(v)$, que não são vizinhos de $v$; logo $v\notin B_{\pi(f(v))}$ e o par é
  viável. O número de cores abertas no caminho é o número de cores de $f$
  mapeadas, no máximo $c^\star$. Assim, as podas não atingem o caminho: a poda
  por $u\ge\UB$ exigiria $c^\star\ge\UB$; a cor nova é oferecida porque
  $u+1\le c^\star<\UB$; e nenhum grupo fica sem opção, pois o par seguido por
  $(S,f)$ continua viável. A busca em profundidade percorre todos os filhos
  não podados, então chega à folha desse caminho, que tem no máximo $c^\star$
  cores, ou antes disso reduz $\UB$ para no máximo $c^\star$ ou para $\LB$.
\end{proof}

\begin{teorema}\label{teo:correcao}
  Sem limite de tempo, o algoritmo termina e devolve uma solução viável com
  $\chiP$ cores. Com limite de tempo, ele devolve uma solução viável com $\UB$
  cores e um valor $\LB$ tais que $\LB\le\chiP\le\UB$.
\end{teorema}

\begin{proof}
  \emph{Viabilidade}: a incumbente inicial vem da heurística e do guloso
  iterado, que só produzem colorações próprias; as folhas da busca também,
  porque um par $(v,c)$ só é usado se $v\notin B_c$. Toda solução ainda é
  conferida pelo validador. Logo $\UB\ge\chiP$ em todo momento.

  \emph{Terminação}: a árvore tem profundidade $k$ (cada nível atende um
  grupo) e cada nó tem no máximo $s\cdot(u+1)$ filhos, então é finita.

  \emph{Limite inferior}: pelos Lemas~\ref{lem:dominancia},
  \ref{lem:clique} e~\ref{lem:cobertura}, o $\LB$ da raiz é um limite
  inferior para $\chiP$. Na estratégia ascendente, $\LB$ só passa a $t+1$
  quando a busca com $\UB=t+1$ termina sem solução; pelo
  Lema~\ref{lem:caminho} com $c^\star\le t$, isso só ocorre se não existe
  solução com no máximo $t$ cores, isto é, se $\chiP>t$.

  \emph{Otimalidade}: na busca descendente sem limite de tempo, suponha que
  ela termine com $\UB>\chiP$. Pelo Lema~\ref{lem:dominancia}, existe solução
  ótima da instância reduzida com $c^\star=\chiP<\UB$ cores, e pelo
  Lema~\ref{lem:caminho} a busca teria encontrado uma solução com no máximo
  $\chiP$ cores ou reduzido $\UB$ para $\LB\le\chiP$, uma contradição. Na
  ascendente, a primeira meta $t$ com solução é ótima, pois $t-1$ já foi
  refutada (ou $t=\LB$). Por fim, a poda por $\UB\le\LB$ só é aplicada
  quando a incumbente já é ótima.
\end{proof}

A correção também foi conferida experimentalmente: os testes automatizados
comparam o resultado do algoritmo com uma força bruta em 300 instâncias
aleatórias pequenas e verificam, em instâncias da literatura com o tempo
esgotando no meio da busca, que $\LB\le\chiP\le\UB$ para o valor ótimo
publicado.

% ---------------------------------------------------------------------------
\section{Análise de complexidade}\label{sec:analise}

Seja $w$ o tamanho da palavra de máquina, de modo que uma operação sobre um
vetor de $n$ bits custa $O(\lceil n/w\rceil)$, e seja $c\le k$ o número de
cores abertas.

\paragraph{Pré-processamento.} Cada rodada da redução por dominância compara
os pares de vértices de cada grupo, com custo
$O(\sum_i|V_i|^2\cdot n/w)\subseteq O(sn^2/w)$, e há no máximo $n$ rodadas,
pois cada rodada que continua remove pelo menos um vértice. A construção de
$H$ custa $O(k^2 s\cdot n/w)$, e a busca da clique de grupos é limitada a
$2\cdot10^5$ nós. A cobertura gulosa por cliques cria no máximo $m$ cliques,
cada uma em $O(n^2/w)$, e a escolha da família custa $O(m^2 s)$ por rodada.
A heurística do Trabalho 1 custa $O(n\log n+cn+m)$, e cada iteração do guloso
iterado, $O(k\,s\,c\cdot n/w)$. Todo o pré-processamento é polinomial.

\paragraph{Custo por nó.} A escolha do grupo percorre os grupos pendentes e,
para cada um, as cores abertas, contando bits: $O(k\,c\cdot n/w)$. A geração
dos filhos custa $O(s\,c)$, e cada filho atualiza um único $B_c$ em
$O(n/w)$. O custo por nó é, portanto, $O(k\,c\cdot n/w)$.

\paragraph{Número de nós.} Cada folha corresponde a uma seleção e a uma
partição dos $k$ grupos em classes de cor rotuladas em ordem canônica (a cor
nova é sempre a de menor índice livre, e as cores da clique $Q$ são fixas).
Assim, há no máximo $\prod_i|V_i|\cdot \mathcal{B}_k$ folhas, onde
$\mathcal{B}_k$ é o número de Bell (o número de partições de um conjunto de
$k$ elementos), e no máximo $k$ vezes isso de nós, pois a profundidade é $k$.
O tempo de pior caso é
\[
  O\Big(\prod_{i=1}^{k}|V_i|\cdot \mathcal{B}_k\cdot k^2 c\cdot \frac{n}{w}\Big),
\]
exponencial em $k$, como esperado para um problema NP-difícil. A quebra de
simetria evita o fator $c^k$ (ou $k!$) de uma enumeração que trata as cores
como distinguíveis. O limite acima é de pior caso: as podas por limite e por
inviabilidade descartam subárvores inteiras, e quando $\LB=\UB$ na raiz nenhum
nó é explorado.

\paragraph{Memória.} A instância usa $O(n+m)$; a busca guarda a atribuição
parcial e os conjuntos $B_c$, em $O(k+c\cdot n/w)$, e a pilha de recursão tem
profundidade $k$. Cada nível guarda a lista de filhos, com $O(s\,c)$ pares,
somando $O(k\,s\,c)$.

% ---------------------------------------------------------------------------
\section{Experimentos}\label{sec:experimentos}

\subsection{Instâncias}

Usamos as mesmas instâncias do primeiro trabalho: o banco clássico do PCP,
distribuído pelo grupo de C. C. Ribeiro e mantido no repositório de
\cite{santini2017}, que acompanha \cite{furini2018}. A
Tabela~\ref{tab:instancias} resume os três conjuntos.

\begin{itemize}
  \item random \cite{frota2010}: grafos aleatórios com grupos de dois
    vértices. Há um conjunto com densidade 0,5 e $n$ de 20 a 120, e outro com
    $n=90$ e densidade de 0,1 a 0,9, com cinco instâncias por combinação.
  \item nsfnet \cite{li2000}: roteamento na rede NSF (14 nós e 21 enlaces).
    Cada conexão pedida com probabilidade $p\in\{0{,}1,\dots,1{,}0\}$ gera um
    grupo com até duas rotas, e duas rotas são adjacentes quando compartilham
    um enlace.
  \item ring \cite{frota2010}: o mesmo processo em redes em anel com 10, 15 e
    20 nós. Os cinco arquivos \texttt{ring\_n\{10,15,20\}p1.0s\{1..5\}} de cada
    tamanho são cópias idênticas, como notam \cite{furini2018}; eles foram
    mantidos para preservar os nomes originais.
\end{itemize}

\begin{table}[ht]
  \centering
  \caption{Conjuntos de instâncias usados nos experimentos.}
  \label{tab:instancias}
  \begin{tabular}{lrcccl}
    \toprule
    Conjunto & Arquivos & $n$ & $k$ & Densidade & Origem \\
    \midrule
    random & 80  & 20--120 & 10--60  & 0,10--0,91 & \cite{frota2010} \\
    nsfnet & 50  & 16--257 & 12--182 & 0,09--0,13 & \cite{li2000} \\
    ring   & 150 & 8--760  & 4--380  & 0,36--0,41 & \cite{frota2010} \\
    \bottomrule
  \end{tabular}
\end{table}

As instâncias chegam a 760 vértices e 117 mil arestas, tamanho suficiente para
que nenhuma enumeração completa seja viável, por isso não foi preciso criar um
gerador. Elas também permitem comparar os resultados com os valores ótimos
publicados. O código, as instâncias e os resultados deste trabalho estão em
\url{https://github.com/Cabetoln/trab2paa}; os do primeiro trabalho, em
\url{https://github.com/pondof/trab1PAA}. As instâncias originais e o código
de \cite{furini2018} estão em
\url{https://github.com/alberto-santini/selective-graph-colouring}.

\subsection{Metodologia}

Os experimentos foram executados em um computador com processador Apple M2
Pro (6 núcleos de desempenho e 4 de eficiência), 16~GB de memória e macOS~26,
usando CPython~3.14.3. O \emph{branch and bound} usou a estratégia híbrida e
um limite de 60~s por instância, contado desde o início do pré-processamento,
de modo que o tempo do guloso iterado entra no total. Para encurtar a
execução, seis instâncias rodaram ao mesmo tempo, cada uma em um processo
próprio com uma única \emph{thread}. A heurística do primeiro trabalho foi
executada na mesma rodada, na mesma máquina. Os tempos foram medidos com
\texttt{time.perf\_counter} e não incluem a leitura do arquivo.

O limite de 60~s foi escolhido para que o experimento completo, com 280
instâncias, possa ser repetido por um leitor em menos de meia hora com seis
processos (ou em cerca de 2,6~h com um só), e é também o padrão da interface
\texttt{main.py}, que qualquer pessoa pode usar para testar uma instância. Esse
valor é baixo perto do que a literatura usa: os métodos exatos de
\cite{furini2018} chegam a uma hora por instância. Os ótimos provados
mostram, portanto, o que o algoritmo resolve em um minuto; como ele é exato e
sempre termina, com mais tempo ele pode provar o ótimo de outras instâncias.
Nas instâncias resolvidas, o maior tempo foi de 27~s, bem abaixo do limite; nas demais, o algoritmo informa o intervalo
$[\LB,\UB]$ em que o ótimo está, e qualquer limite pode ser escolhido com a
opção \texttt{--limite}.

Os dois métodos são determinísticos: o guloso iterado usa um gerador
pseudoaleatório com semente fixa. A única fonte de variação entre execuções é
o limite de tempo, pois uma máquina mais lenta explora menos nós em 60~s e
pode terminar com limites piores. Para reproduzir as tabelas, basta rodar
\texttt{python experimentos.py --limite 60 --processos 6} na pasta do projeto;
o resultado de cada instância vai para \texttt{resultados/bb.csv} e as
tabelas, para \texttt{resultados/tabela\_bb.md}. Todas as soluções das 280
instâncias passaram pelo validador.

A correção da implementação também foi testada com \texttt{python -m unittest
testes}. O teste principal gera 300 instâncias aleatórias pequenas (até 7
grupos de até 3 vértices) e compara o valor do algoritmo com o de uma força
bruta que enumera todas as seleções e colorações. Outro teste interrompe a
busca após 2~s em cinco instâncias difíceis e confere, para as três
estratégias, que $\LB\le\chiP\le\UB$ para o valor ótimo publicado.

Os valores de referência são os mesmos do primeiro trabalho, retirados das
Tabelas~2 e~3 de \cite{furini2018}. Eles cobrem 187 dos 280 arquivos. Em 180
deles o valor é ótimo comprovado; nos outros 7 (\texttt{n90p2t2s4},
\texttt{n90p3t2s4} e as cinco cópias de \texttt{ring\_n20p1.0s1}) é o melhor
valor conhecido. O gap de um método é $100\cdot(c-c^*)/c^*$, onde $c$ é o
número de cores do método e $c^*$ é o valor de referência.

\subsection{Resultados}

A Tabela~\ref{tab:resumo} compara o \emph{branch and bound} (B\&B) com a
heurística do primeiro trabalho (T1) em cada conjunto. As médias de cores e de
tempo são sobre todas as instâncias; os gaps, sobre as instâncias com
referência. A coluna ``Ótimo provado'' conta as instâncias em que o
algoritmo terminou com $\LB=\UB$, e entre parênteses quantas delas foram
resolvidas já na raiz, sem explorar nenhum nó.

\begin{table}[ht]
  \centering
  \caption{Resumo por conjunto: heurística do Trabalho 1 (T1) e
    \emph{branch and bound} (B\&B) com limite de 60~s.}
  \label{tab:resumo}
  \footnotesize\setlength{\tabcolsep}{3.5pt}
  \begin{tabular}{lrrrrrrrrrr}
    \toprule
    & & Ótimo & \multicolumn{3}{c}{Cores (média)} & \multicolumn{2}{c}{Tempo médio} & Com & \multicolumn{2}{c}{Gap médio (\%)} \\
    \cmidrule(lr){4-6}\cmidrule(lr){7-8}\cmidrule(lr){10-11}
    Conjunto & Inst. & provado & T1 & B\&B & $\LB$ & T1 (ms) & B\&B (s) & ref. & T1 & B\&B \\
    \midrule
    nsfnet & 50  & 25 (24) & 9,66  & 8,08  & 7,52  & 0,17 & 30,0 & 32 & 19,17 & 0,00 \\
    random & 80  & 54 (1)  & 9,74  & 7,00  & 6,54  & 0,10 & 23,5 & 56 & 43,30 & 2,37 \\
    ring   & 150 & 48 (40) & 22,70 & 19,73 & 18,08 & 0,71 & 41,1 & 99 & 18,83 & 2,35 \\
    \midrule
    total  & 280 & 127 (65) & & & & & & 187 & & \\
    \bottomrule
  \end{tabular}
\end{table}

A Tabela~\ref{tab:grupos} detalha os resultados por grupo de instâncias. Além
das cores da heurística e do B\&B, ela mostra o limite superior inicial
($\UB_0$, depois do guloso iterado) e o limite inferior final ($\LB$). A
coluna ``Ref.'' traz o número de instâncias do grupo com valor de referência
e, entre parênteses, em quantas delas o B\&B chegou a esse valor. A tabela
completa, instância por instância, está em \texttt{resultados/tabela\_bb.md}.

\begin{table}[ht]
  \centering
  \caption{Resultados por grupo de instâncias: médias de $n$, das cores da
    heurística (T1), do limite superior inicial ($\UB_0$), do B\&B e do
    limite inferior final ($\LB$), e tempo médio do B\&B.}
  \label{tab:grupos}
  \footnotesize\setlength{\tabcolsep}{3.5pt}
  \begin{tabular}{llrrrrrrrrr}
    \toprule
    Conjunto & Grupo & Inst. & $n$ & Ótimo & T1 & $\UB_0$ & B\&B & $\LB$ & Ref. & Tempo (s) \\
    \midrule
    random & $n=20$,  $p=0{,}5$ & 5 & 20  & 5 & 3,4  & 3,0  & 3,0  & 3,0  & 0     & 0,09 \\
           & $n=40$,  $p=0{,}5$ & 5 & 40  & 5 & 6,0  & 4,0  & 4,0  & 4,0  & 0     & 0,28 \\
           & $n=60$,  $p=0{,}5$ & 5 & 60  & 5 & 7,6  & 5,6  & 5,0  & 5,0  & 1 (1) & 0,53 \\
           & $n=70$,  $p=0{,}5$ & 5 & 70  & 5 & 8,4  & 6,0  & 6,0  & 6,0  & 3 (3) & 1,10 \\
           & $n=80$,  $p=0{,}5$ & 5 & 80  & 5 & 9,2  & 7,0  & 6,0  & 6,0  & 5 (5) & 5,63 \\
           & $n=90$,  $p=0{,}5$ & 5 & 90  & 4 & 9,6  & 7,4  & 7,0  & 6,8  & 5 (5) & 28,22 \\
           & $n=100$, $p=0{,}5$ & 5 & 100 & 2 & 10,4 & 8,0  & 7,4  & 6,8  & 5 (3) & 44,59 \\
           & $n=120$, $p=0{,}5$ & 5 & 120 & 0 & 12,0 & 9,0  & 9,0  & 7,0  & 5 (0) & 60,00 \\
           & $n=90$,  $p=0{,}1$ & 5 & 90  & 5 & 3,6  & 3,0  & 2,8  & 2,8  & 5 (5) & 0,59 \\
           & $n=90$,  $p=0{,}2$ & 5 & 90  & 5 & 5,2  & 4,0  & 3,8  & 3,8  & 5 (5) & 1,14 \\
           & $n=90$,  $p=0{,}3$ & 5 & 90  & 5 & 7,0  & 5,0  & 5,0  & 5,0  & 5 (5) & 5,14 \\
           & $n=90$,  $p=0{,}4$ & 5 & 90  & 5 & 8,8  & 6,0  & 5,8  & 5,8  & 5 (5) & 11,84 \\
           & $n=90$,  $p=0{,}6$ & 5 & 90  & 3 & 11,8 & 9,0  & 8,4  & 8,0  & 4 (3) & 36,13 \\
           & $n=90$,  $p=0{,}7$ & 5 & 90  & 0 & 14,0 & 10,8 & 10,0 & 9,0  & 0     & 60,00 \\
           & $n=90$,  $p=0{,}8$ & 5 & 90  & 0 & 17,4 & 13,0 & 12,4 & 11,2 & 4 (2) & 60,00 \\
           & $n=90$,  $p=0{,}9$ & 5 & 90  & 0 & 21,4 & 16,8 & 16,4 & 14,4 & 4 (2) & 60,00 \\
    \midrule
    nsfnet & $p=0{,}1$ & 5 & 26  & 5 & 2,8  & 2,4  & 2,4  & 2,4  & 0     & 0,00 \\
           & $p=0{,}2$ & 5 & 53  & 5 & 4,6  & 4,0  & 4,0  & 4,0  & 0     & 0,13 \\
           & $p=0{,}3$ & 5 & 79  & 4 & 7,0  & 5,4  & 5,4  & 5,2  & 0     & 12,00 \\
           & $p=0{,}4$ & 5 & 103 & 5 & 7,8  & 6,4  & 6,4  & 6,4  & 3 (3) & 0,00 \\
           & $p=0{,}5$ & 5 & 125 & 1 & 9,2  & 7,4  & 7,4  & 6,6  & 4 (4) & 48,00 \\
           & $p=0{,}6$ & 5 & 151 & 2 & 10,6 & 9,0  & 9,0  & 8,4  & 5 (5) & 36,01 \\
           & $p=0{,}7$ & 5 & 181 & 2 & 12,2 & 10,2 & 10,2 & 9,6  & 5 (5) & 36,00 \\
           & $p=0{,}8$ & 5 & 205 & 0 & 13,0 & 11,0 & 11,0 & 9,8  & 5 (5) & 60,00 \\
           & $p=0{,}9$ & 5 & 225 & 1 & 14,2 & 12,0 & 12,0 & 11,2 & 5 (5) & 48,00 \\
           & $p=1{,}0$ & 5 & 249 & 0 & 15,2 & 13,0 & 13,0 & 11,6 & 5 (5) & 60,00 \\
    \midrule
    ring   & 10 nós & 50 & 97  & 35 & 9,32  & 7,98  & 7,98  & 7,62  & 16 (16) & 18,71 \\
           & 15 nós & 50 & 237 & 12 & 21,60 & 18,66 & 18,66 & 17,30 & 38 (23) & 45,65 \\
           & 20 nós & 50 & 421 & 1  & 37,18 & 32,56 & 32,56 & 29,32 & 45 (17) & 58,81 \\
    \bottomrule
  \end{tabular}
\end{table}

Para medir a distância entre a heurística do primeiro trabalho e o ótimo, a
Tabela~\ref{tab:t1otimo} usa todas as instâncias com ótimo conhecido: as 127
em que o B\&B provou o ótimo e as que têm ótimo comprovado na literatura,
246 ao todo. No primeiro trabalho só era possível fazer essa comparação nas
instâncias com referência; o B\&B acrescenta 66 instâncias a esse conjunto.

\begin{table}[ht]
  \centering
  \caption{Número de instâncias com ótimo conhecido segundo a diferença, em
    cores, entre a heurística do Trabalho 1 e o ótimo.}
  \label{tab:t1otimo}
  \small
  \begin{tabular}{lrrrrrr}
    \toprule
    Conjunto & 0 & 1 & 2 & 3 & $\ge4$ & Total \\
    \midrule
    nsfnet & 5  & 16 & 21 & 6  & 0  & 48  \\
    random & 4  & 10 & 21 & 16 & 22 & 73  \\
    ring   & 7  & 25 & 19 & 15 & 59 & 125 \\
    \midrule
    total  & 16 & 51 & 61 & 37 & 81 & 246 \\
    \bottomrule
  \end{tabular}
\end{table}

\subsection{Discussão}

\paragraph{B\&B contra a heurística do Trabalho 1.} O B\&B usou menos cores
que a heurística em 261 das 280 instâncias e empatou nas outras 19. Em média,
a redução foi de 1,6 cor em nsfnet, 2,7 em random e 3,0 em ring, e o gap em
relação à referência caiu de 19\% para 0\% em nsfnet, de 43\% para 2,4\% em
random e de 19\% para 2,4\% em ring. A Tabela~\ref{tab:grupos} mostra de onde
vem essa redução. Em nsfnet e ring, a busca em árvore não melhorou o limite
superior em nenhuma instância: todo o ganho sobre a heurística vem do guloso
iterado ($\UB_0$ igual à coluna B\&B), e a árvore serve para provar que essa
solução é ótima. Em random, a busca melhorou o limite superior inicial em 28
das 80 instâncias. A Tabela~\ref{tab:t1otimo} confirma o que o primeiro
trabalho já indicava: a heurística sozinha atinge o ótimo em só 16 das 246
instâncias com ótimo conhecido, quase todas pequenas, e fica 4 ou mais cores
acima dele em 81, a maioria em ring.

\paragraph{Limites inferiores.} Os dois conjuntos de redes ópticas se
comportam de forma diferente do conjunto random. Em nsfnet e ring, o limite
de cobertura por cliques (que nessas instâncias é o limite de carga dos
enlaces) muitas vezes coincide com o ótimo, e 64 dos 73 ótimos provados
nesses conjuntos saíram já na raiz, quase sempre em menos de 0,1~s; o limite
inferior da raiz só aumentou durante a busca em 10 das 200 instâncias. Nas
instâncias random, os grupos têm dois vértices sem estrutura, e o limite
inferior da raiz fica pequeno (de 2 a 4 cores na maioria dos casos). Lá, 53
dos 54 ótimos foram provados pela busca em árvore, e o limite inferior final
ficou acima do da raiz em 78 das 80 instâncias, quase sempre graças à busca
ascendente, que refuta $t=\LB,\LB+1,\dots$ cores uma a uma. Isso explica por que o tempo cresce rápido com $n$ e com a
densidade em random: com $n=90$ e $p=0{,}5$ já são 28~s em média, e a partir
de $n=120$ ou $p=0{,}7$ nenhuma instância foi resolvida em 60~s.

\paragraph{Comparação com a literatura.} Nas 187 instâncias com referência, o
B\&B chegou ao valor publicado em 132: todas as 32 de nsfnet, 44 de 56 em
random e 56 de 99 em ring. Ele nunca encontrou menos cores que a referência,
e o seu limite inferior nunca passou dela, o que é coerente com a correção
dos dois limites. Quando não chegou, o B\&B ficou 1 cor acima em 32
instâncias, 2 em 14 e de 3 a 5 cores em 9, todas em ring com 20 nós e
$p\ge0{,}9$. O B\&B provou o ótimo de 63 dessas 187 instâncias em até 60~s,
contra 180 com ótimo comprovado nas tabelas de \cite{furini2018}, cujos
métodos exatos levam de segundos a uma hora por instância e usam relaxações
lineares bem mais fortes que os nossos limites combinatórios. Em dois casos o
B\&B foi além da referência: para \texttt{n90p2t2s4} e \texttt{n90p3t2s4}, as
tabelas trazem 4 e 5 cores apenas como melhor valor conhecido, e o B\&B prova
que esses valores são ótimos, em 2,0~s e 9,8~s. Ele também provou o ótimo de
64 das 93 instâncias que não aparecem nas tabelas do artigo.

\paragraph{Tempo.} A heurística do primeiro trabalho levou no máximo 3,8~ms
por instância e 0,12~s para as 280. O B\&B é várias ordens de grandeza mais
lento: somando todas as instâncias, foram 2,6~h de processamento, ou 27~min
de tempo real com seis instâncias em paralelo. Esse tempo se concentra nas
153 instâncias em que o ótimo não foi provado, que usam os 60~s inteiros e
respondem por 2,55~h do total. Nas 127 com ótimo provado, a mediana do tempo foi de
0,002~s em nsfnet, 0,016~s em ring e 1,05~s em random, e o maior tempo foi de
27~s. Entre as instâncias resolvidas está \texttt{ring\_n15p0.6s3}, com 258
vértices, provada em 0,07~s. Nas 153 instâncias não resolvidas, o B\&B termina
com uma solução e um intervalo $[\LB,\UB]$; a diferença média é de 1,1 cor em
nsfnet, 1,4 em random e 2,4 em ring.

% ---------------------------------------------------------------------------
\section{Conclusão}\label{sec:conclusao}

Este trabalho apresentou um algoritmo exato do tipo \emph{branch and bound}
para o Problema de Coloração Seletiva de Grafos. O algoritmo adapta ao
problema a ramificação do DSATUR, quebra a simetria entre cores com uma
clique de grupos, descarta vértices dominados e combina dois limites
inferiores combinatórios com um limite superior obtido pela heurística do
primeiro trabalho e pelo guloso iterado. Provamos que ele devolve o ótimo
quando termina e que, interrompido pelo limite de tempo, devolve uma solução
viável e um limite inferior válido. A implementação usa apenas a biblioteca
padrão do Python.

Nas 280 instâncias da literatura, com 60~s por instância, o algoritmo provou
o ótimo em 127 e chegou ao valor publicado em 132 das 187 instâncias com
referência, incluindo todas as de nsfnet. Ele também provou a otimalidade de
dois valores que a literatura trazia só como melhores conhecidos e de 64
instâncias que não aparecem nas tabelas de \cite{furini2018}. Comparado com a
heurística do primeiro trabalho, usou menos cores em 261 instâncias, e o gap
médio caiu de 19\% (nsfnet e ring) e 43\% (random) para no máximo 2,4\%. O
custo é o tempo: a heurística resolve todas as instâncias em 0,12~s, enquanto
o B\&B usa os 60~s inteiros em 153 delas.

As instâncias não resolvidas são as random com densidade a partir de 0,7 ou
com $n\ge100$, nas quais
os limites inferiores combinatórios ficam longe do ótimo, e as ring maiores,
em que a árvore é grande demais para 60~s em Python. Como trabalhos futuros,
sugerimos: (i) recalcular o limite de cobertura por cliques em cada nó,
considerando só os grupos ainda pendentes, para podar mais cedo nas
instâncias de redes ópticas; e (ii) usar a relaxação linear da formulação de
cobertura, como em \cite{furini2018}, como limite inferior nas instâncias
random, o que exigiria implementar um resolvedor de programação linear sem
bibliotecas externas.

\bibliographystyle{sbc}
\bibliography{referencias}

\end{document}
